\SourceLeaf{110}
\Article{18}
Stellt man die zweite Nadel an verschiedenen Punkten derselben Geraden auf, sodass $\psi$ und $U$ unverändert bleiben und sich nur der Abstand $R$ ändert, und beobachtet man die Ablenkungen der beweglichen Nadel aus ihrer Gleichgewichtslage bei Abwesenheit der zweiten Nadel, also die Winkel $u-u^0$, so lassen sich daraus offenbar durch Elimination die Werte der Koeffizienten $F,F',F''$ usw. ermitteln, soweit sie noch merklich ins Gewicht fallen. Damit erhalten wir
\[
\frac MT=\left(1+\frac{\theta}{Tm}\right)\frac{F}{n\cos(\psi-U)\sin(\psi-u^0)+\sin(\psi-U)\cos(\psi-u^0)}
\]
Dabei lässt sich der Wert von $\frac{\theta}{Tm}$ nach dem in Artikel $8$ beschriebenen Verfahren bestimmen. Für eine einfachere praktische Durchführung lohnt es sich, Folgendes zu beachten.

I. Anstelle eines Vergleichs von $u$ mit $u^0$ empfiehlt es sich, zwei entgegengesetzte Ablenkungen miteinander zu vergleichen. Dazu kehrt man die Lage der zweiten Nadel um, lässt also $R$ und $\psi$ unverändert und vergrößert den Winkel $U$ um zwei rechte Winkel. Bezeichnen wir die zu diesen Lagen gehörenden Werte von $u$ mit $u',u''$, so wäre bei vollkommener Symmetrie \textit{genau} $u''=-u'$, sofern zugleich $u^0=0$ wäre. Es ist aber überflüssig, diese Bedingungen peinlich genau einzuhalten. Denn $u'$ und $u''$ werden offenbar durch ähnliche Reihen bestimmt, deren erste Terme \textit{genau} entgegengesetzte Werte haben. Damit werden auch $\tfrac12(u'-u'')$ und $\operatorname{tang}\tfrac12(u'-u'')$ durch eine ähnliche Reihe bestimmt, in der der Koeffizient des ersten Terms \textit{genau} $=F$ ist.

II. Noch besser ist es, stets vier Versuche zusammenzufassen und dabei auch den Winkel $\psi$ um zwei rechte Winkel zu ändern oder den Abstand $R$ auf der anderen Seite zu nehmen. Entsprechen den beiden letzten Versuchen die Werte von $u$, nämlich $u''',u''''$, so wird auch die Differenz $\tfrac12(u'''-u'''')$ durch eine ähnliche Reihe ausgedrückt, deren erster Term ebenfalls den Koeffizienten $=F$ hat. Man beachte, wie sich aus dem Vorangehenden leicht ergibt, dass für ungerades $n$ die Koeffizienten $F,F'',F''''$ usw. in den beiden Reihen für $u'-u^0$ und $u'''-u^0$ bis ins Unendliche genau gleich und die Koeffizienten $F',F''',F^{\mathrm{v}}$ usw. bis ins Unendliche genau entgegengesetzt wären. Entsprechendes würde für $u''-u^0$ und $u''''-u^0$ gelten, sodass in der Reihe für $u'-u''+u'''-u''''$ jeder zweite Term wegfiele. In der Natur, wo $n=2$ ist, gilt diese Beziehung zwischen den Reihen für $u'-u^0$ und $u'''-u^0$ im Allgemeinen jedoch nicht genau. Schon bei der Potenz $R^{-6}$ ergeben sich Koeffizienten, die nicht genau entgegengesetzt sind. Man kann jedoch zeigen, dass auch dieser Term in der Kombination $u'-u''+u'''-u''''$ vollständig wegfällt, sodass $\operatorname{tang}\tfrac14(u'-u''+u'''-u'''')$ die Form hat
\[
LR^{-3}+L'R^{-5}+L''R^{-7}+\text{usw.}
\]
\SourceLeaf{111}oder allgemeiner, wenn wir den Wert von $n$ vorerst offenlassen, die Form
\[
LR^{-(n+1)}+L'R^{-(n+3)}+L''R^{-(n+5)}+\text{usw.}
\]
mit $L=F$.

III. Die Winkel $\psi,U$ wählt man zweckmäßig so, dass kleine Fehler bei ihrer Messung den Wert von $F$ nicht merklich ändern. Dazu muss der Wert von $U$ für einen \textit{gegebenen} Wert von $\psi$ so gewählt werden, dass $F$ ein Maximum annimmt. Es muss also gelten
\[
\operatorname{cotang}(\psi-U)=n\operatorname{tang}(\psi-u^0)
\]
Daraus ergibt sich
\[
F=\pm\frac{mM}{mT+\theta}\BareRoot{(nn\sin(\psi-u^0)^2+\cos(\psi-u^0)^2)}
\]

Der Winkel $\psi$ wiederum muss so gewählt werden, dass dieser Wert von $F$ entweder ein Maximum oder ein Minimum annimmt. Das Maximum tritt bei $\psi-u^0=90^0$ oder $270^0$ ein, wo gilt
\[
F=\pm\frac{nmT}{mT+\theta}
\]
und das Minimum bei $\psi-u^0=0$ oder $180^0$, wo gilt
\[
F=\pm\frac{mM}{mT+\theta}
\]

